\hnsection{自由李代数}

最后还需谈及一系列关于李代数的工作，其中李代数与李群理论的联系十分微弱；另一方面，这些研究对于“抽象”群理论，尤其是幂零群理论，有着重要应用。

其源头是 P. Hall {[}24{]} 的工作，该工作发表于 1932 年。不过这里并未讨论李代数：P. Hall 想研究一类 p-群，即他称为“正则”的那些群。但这使他详细考察了群的迭代换位子和下中心列；关于此事，他建立了 Jacobi 恒等式的一个版本（cf.~第二章，§ 4, no. 4, 公式 (20)）以及“Hall 公式”

\[
(x y) ^ {n} = x ^ {n} y ^ {n} (x, y) ^ {n (1 - n) / 2} \dots \quad (\text { cf.   Chapter   II }, \S 5, \text { Exercise   9 }).
\]

几乎与此同时（1935--1937 年），W. Magnus ({[}25 a{]} 和 {[}25 b{]}) 以及 E. Witt {[}30{]} 的奠基性工作相继出现。在 {[}25 a{]} 中，Magnus 使用了与 Hausdorff 相同的形式幂级数代数 \(\hat{\mathbf{A}}\)（后来称为“Magnus 代数”）；他把自由群 F 嵌入其中，并利用 \(\hat{\mathbf{A}}\) 的自然滤过得到 F 的子群的一个递减序列 \((\mathbf{F}_{n})\)；这是滤过的最早例子之一。他猜想 \(F_{n}\) 与 F 的下中心列各项重合。这一猜想在他的第二篇论文中得到证明 {[}25 b{]}；在这项工作中，他还明确展示了自己的思想与 P. Hall 的思想之间的密切关系，并定义了自由李代数 L（作为 \(\hat{\mathbf{A}}\) 的子代数），他基本上证明了该代数可与 F 的分次代数认同。在 {[}30{]} 中，Witt 在多个方面完成了这一结果。特别是，他证明 L 的包络代数是自由结合代数，并立即推出 L 的齐次分量的秩（“Witt 公式”）。

至于 L 的基的证明，即所谓“Hall 基”（cf.~第二章，§ 2, no. 11），似乎直到 1950 年才出现在 M. Hall 的一篇短文中 {[}40{]}，尽管它已经隐含在上面引用的 P. Hall 和 W. Magnus 的工作中。

